\documentclass[10pt,a4paper]{amsart}
\usepackage[margin=19mm]{geometry}
\usepackage{amsmath,amssymb,mathtools}
\usepackage[scr=rsfs,cal=euler]{mathalfa}
\usepackage{xfrac}
\usepackage[unicode,hidelinks,bookmarks=false]{hyperref}
\newcommand{\ie}{i.e.\ }
\newcommand{\cf}{cf.\ }
\newcommand{\resp}{respectively\ }
\newcommand{\Subsection}{\subsection}
\providecommand{\nobreakdash}{\nobreak\hbox{-}}
\setlength{\emergencystretch}{2em}
\multlinegap0pt
\newtheorem{proofof}{Proofof}
\usepackage{mathtools}
\usepackage{amssymb}
\usepackage{xcolor}
\usepackage{csquotes}
\usepackage[shortlabels]{enumitem}
\newtheorem{lemma}{Lemma}
\renewcommand{\geq}{\geqslant}
\renewcommand{\leq}{\leqslant}
\newcommand{\qgauss}[2]{\genfrac{[}{]}{0pt}{}{#1}{#2}_q}
\title{Subgroup bounds for abelian $p$-groups with applications}
\author{Stefanos Aivazidis}
\date{Source: arXiv:2606.03320v1 (CC BY 4.0)}
\begin{document}
\maketitle

\noindent\emph{Source and licence.} Subgroup bounds for abelian $p$-groups with applications. Version arXiv:2606.03320v1 (CC BY 4.0), distributed by arXiv under the Creative Commons Attribution 4.0 International licence, http://creativecommons.org/licenses/by/4.0/, as recorded in the arXiv licence metadata for this exact version and shown on the abstract page https://arxiv.org/abs/2606.03320v1. Full licence notice: \texttt{licenses/arxiv-2606.03320-CC-BY-4.0.txt}.\par\smallskip

\noindent\textit{Editorial scope.} Counted unit: the article's lemma lem:shifted-chain (source0.tex lines 667--685). The article's complete original proof of it is delivered with it (source0.tex lines 687--782, one \texttt{proof} environment of 96 lines). No mathematics is written, repaired or re-derived: the statement and the proof are the article's own lines, in the article's own notation, and no retained line is substituted or re-flowed.  Retained uncounted as the source-local context the proof needs: lines 622--628.  Nothing was omitted from any retained span, so the package carries no omission record.  No retained line contains a citation, so the wrapper carries no bibliography and no bibliography entry.  No retained line contains a figure environment, an image-inclusion command or drawing code, so no image is delivered and the package declares no asset dependency.  Editorial formatting only, in addition: an AMS \texttt{amsart} preamble that restates the article's own declarations and macros used by the retained lines, namely the environments \texttt{lemma} on separate counters, together with the standard packages those lines need.  The retained environments print in this document as Lemma 1 in order of appearance, whereas the article prints the same items with the numbering of its own full document.  The complete native source of the article as distributed in the arXiv e-print for this exact version is bundled unchanged as \texttt{sources/201963/source0.tex}.
\begin{equation}\label{eq:durfee-chain-shifted-square-sum}
\sum_{\substack{t\geq 0\\ a\geq y_1\geq\cdots\geq y_t\geq 1}}
        q^{\sum_{j=1}^t y_j(y_j+1)}
        \mathcal D_a(\mathbf y)
        =
        \prod_{r=2}^{a+1}\frac{1}{1-q^r}.
\end{equation}

\begin{lemma}\label{lem:shifted-chain}
For $a\geq 1$ and every integer $N\geq 0$,
\[
\sum_{\substack{0\leq t\leq N\\ a\geq y_1\geq\cdots\geq y_t\geq 1}}
        q^{\sum_{j=1}^t y_j(y_j-1)}
        \mathcal D_a(\mathbf y)
        \leq
        \frac{N+1}{(q;q)_{a-1}}.
\]
Moreover,
\[
\lim_{N\to\infty}\frac{1}{N+1}
\sum_{\substack{0\leq t\leq N\\ a\geq y_1\geq\cdots\geq y_t\geq 1}}
        q^{\sum_{j=1}^t y_j(y_j-1)}
        \mathcal D_a(\mathbf y)
        =
        \frac{1}{(q;q)_{a-1}}.
\]
\end{lemma}

\begin{proofof}
Let $L_{a,N}^{(c)}$ be the part of the left side contributed by chains of
length at most $N$ with exactly $c$ terminal entries equal to $1$. Suppose
first that $c\geq 1$. Delete these terminal ones and subtract $1$ from every
remaining entry. What remains is a chain of length $s\leq N-c$,
\[
        a-1\geq z_1\geq\cdots\geq z_s\geq 1.
\]
The deleted terminal ones have zero shifted weight, since $1(1-1)=0$, and the
remaining shifted weight becomes
\[
        \sum_{i=1}^s (z_i+1)z_i=
        \sum_{i=1}^s z_i(z_i+1).
\]

We also need the Gaussian product after the same deletion. If $s\geq 1$, then
expanding the Gaussian coefficients in $q$-Pochhammer symbols gives
\[
\begin{aligned}
        \mathcal D_a(z_1+1,\ldots,z_s+1,\underbrace{1,\ldots,1}_{\text{$c$ entries}})
        &=
        \qgauss{a}{z_1+1}
        \prod_{i=2}^{s}\qgauss{z_{i-1}+1}{z_i+1}
        \qgauss{z_s+1}{1}                                      \\
        &=
        \frac{(q;q)_a}
        {(q;q)_1(q;q)_{a-1-z_1}
        \prod_{i=2}^{s}(q;q)_{z_{i-1}-z_i}(q;q)_{z_s}}           \\
        &=
        \qgauss{a}{1}\mathcal D_{a-1}(z_1,\ldots,z_s).
\end{aligned}
\]
The same identity holds when $s=0$, in which case both sides are
$\qgauss{a}{1}$. Thus, for fixed $c\geq 1$, the contribution is at most
\[
        \qgauss{a}{1}
        \sum_{\substack{t\geq 0\\ a-1\geq z_1\geq\cdots\geq z_t\geq 1}}
        q^{\sum_i z_i(z_i+1)}\mathcal D_{a-1}(z_1,\ldots,z_t),
\]
which, by \eqref{eq:durfee-chain-shifted-square-sum} with $a-1$ in place of $a$, is
\[
        \qgauss{a}{1}
        \prod_{r=2}^{a}\frac{1}{1-q^r}
        =
        \frac{1-q^a}{1-q}
        \prod_{r=2}^{a}\frac{1}{1-q^r}
        =
        \frac{1}{(q;q)_{a-1}}.
\]

It remains, for the inequality, to treat the chains with no terminal entry equal
to $1$. Appending one terminal $1$ to such a chain is an injective map into
the chains with exactly one terminal $1$, and the shifted exponent is
unchanged. If the original chain is empty, the Gaussian product is multiplied by
$\qgauss{a}{1}$; otherwise it is multiplied by $\qgauss{y_t}{1}$, where
$y_t\geq 2$ is the last entry of the original chain. In both cases the
multiplier is at least $1$, whence the $c=0$ contribution is bounded by the
unrestricted contribution of chains with one terminal $1$, which is bounded by
the quantity just computed. Since $0\leq c\leq N$, the asserted inequality
follows.

We turn to the limiting assertion. Put
\[
        R_M\coloneqq
        \sum_{\substack{0\leq s\leq M\\ a-1\geq z_1\geq\cdots\geq z_s\geq 1}}
        q^{\sum_i z_i(z_i+1)}\mathcal D_{a-1}(z_1,\ldots,z_s).
\]
By \eqref{eq:durfee-chain-shifted-square-sum} with $a-1$ in place of $a$,
the increasing sequence $R_M$ tends to
\[
        L=\prod_{r=2}^a(1-q^r)^{-1}.
\]
For the limiting assertion we keep the length restriction instead of discarding
it. Then, for $c\geq 1$, the deletion map above is exact, and therefore
$L_{a,N}^{(c)}=\qgauss{a}{1}R_{N-c}$. Hence
\[
        \sum_{c=1}^N L_{a,N}^{(c)}
        =
        \qgauss{a}{1}\sum_{M=0}^{N-1}R_M.
\]
The $c=0$ contribution is bounded independently of $N$, and so vanishes
after division by $N+1$. 
Since $R_M\to L$, the averages of the $R_M$'s also tend to $L$, and hence
\[
        \frac{1}{N+1}\sum_{M=0}^{N-1}R_M \longrightarrow L.
\]
Multiplying by $\qgauss{a}{1}$ gives
\[
        \qgauss{a}{1}L
        =
        \frac{1-q^a}{1-q}\prod_{r=2}^a\frac{1}{1-q^r}
        =
        \frac{1}{(q;q)_{a-1}},
\]
which is the stated limit.
\end{proofof}

\end{document}
